% from ega4-20-fr.tex
\oldpage[IV-4]{255}
En effet, l'hypothèse entraîne, en vertu de \sref{IV.11.3.10}, que pour tout $s\in S$, $f_s:X'_s\to X_s$ est plat ; donc, pour toute section $t\in T_{X/S}(U)$, si $t'$ est son image réciproque dans $\Gamma(f^{-1}(U),\sh{O}_{X'})$, $t'_{s}$, qui est l'image réciproque de $t_s$, est une section \emph{régulière} de $\sh{O}_{X'_s}$ au-dessus de $f^{-1}(U)\cap X'_s$, en vertu de \sref{IV.20.1.12} ; on en conclut que par définition $t'\in T_{X'/S}(f^{-1}(U))$, d'où la proposition.
\end{proof}

On déduit de là, comme dans \sref{IV.20.1.12}, un homomorphisme canonique de $\sh{O}_{X'}$-Algèbres
\[
\label{IV.20.6.8.1}
  f^*(\sh{M}_{X/S})\to\sh{M}_{X'/S}.
  \tag{20.6.8.1}
\]

\begin{env}[20.6.9]
\label{IV.20.6.9}
Considérons enfin un morphisme \emph{quelconque} $S'\to S$, et posons $X'=X\times_S S'$, qui est plat et localement de présentation finie sur $S'$ ; soit $p:X'\to X$ la projection canonique.
Soient $U$ un ouvert de $X$, $t$ une section appartenant à $T_{X/S}(U)$, $t'$ son image dans $\Gamma(p^{-1}(U),\sh{O}_{X'})$ ; pour tout $s'\in S'$, si $s\in S$ est l'image de $s'$, on a $X'_{s'}=X_s\otimes_{k(s)}k(s')$, donc le morphisme $X'_{s'}\to X_s$ est plat, et par suite \sref{IV.20.1.12} l'image réciproque $t'_{s'}$ de $t_s$ dans $\Gamma(p^{-1}(U)\cap X'_{s'},\sh{O}_{X'_{s'}})$ est régulière ; ceci prouve que l'on a $t'\in T_{X'/S'}(p^{-1}(U))$.
Ceci permet de définir canoniquement, comme dans \sref{IV.20.6.8}, un homomorphisme canonique de $\sh{O}_{X'}$-Algèbres
\[
\label{IV.20.6.9.1}
  p^*(\sh{M}_{X/S})\to\sh{M}_{X'/S'}.
  \tag{20.6.9.1}
\]

Moyennant l'identification permise par \sref{IV.20.6.3}, cette notion de changement de base pour les fonctions méromorphes relatives est un cas particulier de la notion analogue pour les pseudo-morphismes relatifs \sref{IV.20.5.8}.
\end{env}

% (continued in ega4-21a-fr.tex, before its first page marker)
\section{Diviseurs}
\label{section:IV.21}

Sur le contenu du présent paragraphe, voir les commentaires de l'introduction du \textsection~20.
Pour les propriétés \emph{globales} des diviseurs, le lecteur se reportera au paragraphe qui leur est consacré au chap.~V.

\subsection{Diviseurs sur un espace annelé}
\EGAReaderTarget{ega4.c04.subsec.21-1}{subsection:IV.21.1}
\label{IV.21.1}

\begin{env}[21.1.1]
\label{IV.21.1.1}
Soient $(X,\sh{O}_X)$ un espace annelé, $\sh{M}_X$ le faisceau des germes de fonctions méromorphes sur $X$ \sref{IV.20.1.3}, $\sh{M}_X^*$ le faisceau (de groupes multiplicatifs) des germes de fonctions méromorphes régulières sur $X$ \sref{IV.20.1.8}.
Il est clair que le faisceau (de groupes multiplicatifs) $\sh{O}_X^*$ des germes de sections inversibles de $\sh{O}_X$ s'identifie canoniquement à un sous-faisceau (de groupes multiplicatifs) de $\sh{M}_X^*$.
\end{env}

\begin{definition}[21.1.2]
\EGAReaderTarget{ega4.c04.statement.21-1-2}{IV.21.1.2}
On appelle \emph{faisceau des diviseurs sur $X$} et on note $\sh{D}iv_X$ le faisceau quotient (de groupes commutatifs) $\sh{M}_X^*/\sh{O}_X^*$;
les sections de ce faisceau au-dessus de $X$ se nomment les \emph{diviseurs sur $X$}; ils forment un groupe commutatif noté $\operatorname{Div}(X)$.
Pour toute section $f$ de $\sh{M}_X^*$ au-dessus de $X$ (autrement dit, toute fonction méromorphe régulière sur $X$ \sref{IV.20.1.8}) on appelle \emph{diviseur de $f$} et on note $\operatorname{div}(f)$ (ou $\operatorname{div}_X(f)$) le diviseur sur $X$ image de $f$ par l'homomorphisme canonique
\[
  \Gamma(X,\sh{M}_X^*)\longrightarrow\Gamma(X,\sh{D}iv_X).
\]
\end{definition}

